% Related lecture: SC2005-L06
% Source: Part 3 (CPU Scheduling), pp. 15--16, 19, 28--29
% Human verified: Yes

\exercisemeta
  {SC2005-L06-EX}
  {2026-08-28}
  {Part 3 (CPU Scheduling), pp.~15--16、19、28--29，lecture examples}
  {题干、Gantt charts 与答案均已核对}
  {已确认（2026-08-28）}

\exercisequestion{SC2005-L06-E01}{Single-core SJF 与 SRTF}

\textbf{对应知识点：}\relatedtopic{SC2005-L06-T01}{SJF 与 SRTF}

\subsubsection*{题目（Lecture Example）}

给定 single-core processes：
\[
\begin{array}{c|cccc}
  \text{Process} & P_1 & P_2 & P_3 & P_4\\ \hline
  A_i & 0 & 2 & 4 & 5\\
  B_i & 7 & 4 & 1 & 4
\end{array}
\]
分别使用 nonpreemptive SJF 与 SRTF，画出 Gantt charts 并计算 average waiting time。

\begin{checkedsolution}
\begin{center}
  \resizebox{0.96\textwidth}{!}{\begin{tikzpicture}[x=0.63cm,y=0.95cm,>=Latex,font=\small]
  \node[anchor=east,font=\bfseries] at (-0.25,1.45) {SJF};
  \draw[fill=blue!12] (0,1.05) rectangle (7,1.85) node[midway] {$P_1$};
  \draw[fill=green!15] (7,1.05) rectangle (8,1.85) node[midway,font=\scriptsize] {$P_3$};
  \draw[fill=orange!18] (8,1.05) rectangle (12,1.85) node[midway] {$P_2$};
  \draw[fill=red!12] (12,1.05) rectangle (16,1.85) node[midway] {$P_4$};
  \foreach \t in {0,7,8,12,16}{\draw (\t,1.05)--(\t,0.92) node[below,font=\scriptsize]{\t};}

  \node[anchor=east,font=\bfseries] at (-0.25,0.05) {SRTF};
  \draw[fill=blue!12] (0,-0.35) rectangle (2,0.45) node[midway] {$P_1$};
  \draw[fill=orange!18] (2,-0.35) rectangle (4,0.45) node[midway] {$P_2$};
  \draw[fill=green!15] (4,-0.35) rectangle (5,0.45) node[midway,font=\scriptsize] {$P_3$};
  \draw[fill=orange!18] (5,-0.35) rectangle (7,0.45) node[midway] {$P_2$};
  \draw[fill=red!12] (7,-0.35) rectangle (11,0.45) node[midway] {$P_4$};
  \draw[fill=blue!12] (11,-0.35) rectangle (16,0.45) node[midway] {$P_1$};
  \foreach \t in {0,2,4,5,7,11,16}{\draw (\t,-0.35)--(\t,-0.48) node[below,font=\scriptsize]{\t};}
\end{tikzpicture}
}

  \smallskip
  \textit{同一 workload 下，nonpreemptive SJF 与 SRTF 的执行顺序。}
\end{center}

Nonpreemptive SJF：$P_1$ 在 $t=0$ 是唯一 Ready process，因此必须先运行；之后选择 $P_3,P_2,P_4$（等长时按 arrival order）。
\[
  (W_1,W_2,W_3,W_4)=(0,6,3,7),\qquad \overline W=4.
\]

SRTF 在 $t=2$ 以 $P_2(4)<P_1(5)$ 抢占，在 $t=4$ 又以 $P_3(1)<P_2(2)$ 抢占。由 $W_i=C_i-A_i-B_i$：
\[
  (W_1,W_2,W_3,W_4)=(9,1,0,2),\qquad \overline W=3.
\]
\end{checkedsolution}

\textbf{检查点：}尚未到达的 process 没有资格参与 SJF 选择；SRTF 比较的是 remaining time，而不是 original burst time。

\exercisequestion{SC2005-L06-E02}{Single-core RR，q=20}

\textbf{对应知识点：}\relatedtopic{SC2005-L06-T03}{Round Robin 与 Time Quantum}

\subsubsection*{题目（Lecture Example）}

$P_1,P_2,P_3,P_4$ 均在 $t=0$ 按此顺序进入 ready queue，burst lengths 为 $(53,17,68,24)$。使用 RR 且 $q=20$，画出 schedule，计算 completion、waiting 与 response times。

\begin{checkedsolution}
\begin{center}
  \resizebox{0.96\textwidth}{!}{\begin{tikzpicture}[x=0.13cm,y=0.92cm,>=Latex,font=\small]
  \node[anchor=east,font=\bfseries] at (-2,1.45) {$0$--$77$};
  \draw[fill=blue!12] (0,1.05) rectangle (20,1.85) node[midway] {$P_1$};
  \draw[fill=orange!18] (20,1.05) rectangle (37,1.85) node[midway] {$P_2$};
  \draw[fill=green!15] (37,1.05) rectangle (57,1.85) node[midway] {$P_3$};
  \draw[fill=red!12] (57,1.05) rectangle (77,1.85) node[midway] {$P_4$};
  \foreach \t in {0,20,37,57,77}{\draw (\t,1.05)--(\t,0.92) node[below,font=\scriptsize]{\t};}

  \node[anchor=east,font=\bfseries] at (-2,0.05) {$77$--$162$};
  \draw[fill=blue!12] (0,-0.35) rectangle (20,0.45) node[midway] {$P_1$};
  \draw[fill=green!15] (20,-0.35) rectangle (40,0.45) node[midway] {$P_3$};
  \draw[fill=red!12] (40,-0.35) rectangle (44,0.45) node[midway,font=\tiny] {$P_4$};
  \draw[fill=blue!12] (44,-0.35) rectangle (57,0.45) node[midway] {$P_1$};
  \draw[fill=green!15] (57,-0.35) rectangle (77,0.45) node[midway] {$P_3$};
  \draw[fill=green!24] (77,-0.35) rectangle (85,0.45) node[midway,font=\scriptsize] {$P_3$};
  \foreach \x/\t in {0/77,20/97,40/117,44/121,57/134,77/154,85/162}{\draw (\x,-0.35)--(\x,-0.48) node[below,font=\scriptsize]{\t};}
\end{tikzpicture}
}

  \smallskip
  \textit{RR schedule 分成两行显示；第二行的时间从 $77$ 延续。}
\end{center}

\[
  (C_1,C_2,C_3,C_4)=(134,37,162,121),
\]
\[
  (W_1,W_2,W_3,W_4)=C_i-B_i=(81,20,94,97),\qquad \overline W=73.
\]
第一次获得 CPU 的时间为 $(0,20,37,57)$，所以 average response time 为 $28.5$。$P_3$ 在 $134$--$154$ 后只剩 $8$；$t=154$ 时没有其他 Ready process，故 scheduler 再次选择 $P_3$，两段可连续绘制，但 quantum boundary 仍存在。
\end{checkedsolution}

\textbf{检查点：}RR 按 ready-queue order 轮换，不按 remaining time 排序；waiting time 可用 $C_i-A_i-B_i$ 自检。

\exercisequestion{SC2005-L06-E03}{Dual-core SRTF 与 RR}

\textbf{对应知识点：}\relatedtopic{SC2005-L06-T05}{Partitioned 与 Global Multiprocessor Scheduling}

\subsubsection*{题目（Lecture Examples）}

在 global dual-core model 下：
\begin{enumerate}
  \item 对 E01 的 arrival/burst data 使用 SRTF；
  \item 对 E02 的四个同时到达 processes 使用 RR，$q=20$。
\end{enumerate}
画出两个 cores 的 schedules，并说明调度事件如何发生。

\begin{checkedsolution}
\begin{center}
  \resizebox{0.96\textwidth}{!}{\begin{tikzpicture}[x=0.86cm,y=0.78cm,>=Latex,font=\small]
  \node[anchor=west,font=\bfseries] at (0,4.25) {Dual-core SRTF};
  \node[anchor=east] at (-0.2,3.35) {Core 1};
  \draw[fill=blue!12] (0,3.0) rectangle (4,3.7) node[midway] {$P_1$};
  \draw[fill=green!15] (4,3.0) rectangle (5,3.7) node[midway,font=\scriptsize] {$P_3$};
  \draw[fill=blue!12] (5,3.0) rectangle (8,3.7) node[midway] {$P_1$};
  \node[anchor=east] at (-0.2,2.35) {Core 2};
  \draw (0,2.0) rectangle (2,2.7);
  \draw[fill=orange!18] (2,2.0) rectangle (6,2.7) node[midway] {$P_2$};
  \draw[fill=red!12] (6,2.0) rectangle (10,2.7) node[midway] {$P_4$};
  \foreach \t in {0,2,4,5,6,8,10}{\draw (\t,2.0)--(\t,1.88) node[below,font=\scriptsize]{\t};}

  \node[anchor=west,font=\bfseries] at (0,0.95) {Dual-core RR, $q=20$};
  \node[anchor=east] at (-0.2,0.1) {Core 1};
  \draw[fill=blue!12] (0,-0.25) rectangle (2.0,0.45) node[midway] {$P_1$};
  \draw[fill=red!12] (2.0,-0.25) rectangle (4.0,0.45) node[midway] {$P_4$};
  \draw[fill=green!15] (4.0,-0.25) rectangle (6.0,0.45) node[midway] {$P_3$};
  \draw[fill=blue!12] (6.0,-0.25) rectangle (7.3,0.45) node[midway] {$P_1$};
  \node[anchor=east] at (-0.2,-0.9) {Core 2};
  \draw[fill=orange!18] (0,-1.25) rectangle (1.7,-0.55) node[midway] {$P_2$};
  \draw[fill=green!15] (1.7,-1.25) rectangle (3.7,-0.55) node[midway] {$P_3$};
  \draw[fill=blue!12] (3.7,-1.25) rectangle (5.7,-0.55) node[midway] {$P_1$};
  \draw[fill=red!12] (5.7,-1.25) rectangle (6.1,-0.55) node[midway,font=\tiny] {$P_4$};
  \draw[fill=green!15] (6.1,-1.25) rectangle (8.9,-0.55) node[midway] {$P_3$};
  \node[anchor=west,font=\scriptsize] at (9.25,-0.15) {Core 1: $0,20,40,60,73$};
  \node[anchor=west,font=\scriptsize] at (9.25,-0.95) {Core 2: $0,17,37,57,61,89$};
\end{tikzpicture}
}

  \smallskip
  \textit{Global dual-core SRTF 与 RR。空白区间表示该 core idle。}
\end{center}

SRTF 在 $t=4$ 从 $(P_1:3,P_2:2,P_3:1)$ 中选择 $P_2,P_3$，令 $P_1$ 等待；$t=5$ 时选择 $(P_2:1,P_1:3)$，令新到达的 $P_4:4$ 等待。最终
\[
  (W_1,W_2,W_3,W_4)=(1,0,0,1),\qquad \overline W=0.5.
\]

Dual-core RR 的 completion times 为
\[
  (C_1,C_2,C_3,C_4)=(73,17,89,61),
\]
故 waiting times 为 $(20,0,21,37)$，average waiting time 为 $19.5$。两个 cores 在各自 idle 的时刻独立触发 scheduler，并非以同步 rounds 推进；同一 process 不能同时占用两个 cores。
\end{checkedsolution}

\textbf{检查点：}每个 arrival、completion 或 quantum expiry 都可能改变 global ready queue；先按真实事件时间更新队列，再给 idle core 分派任务。
